\chapter{Homological Algebra}

\begin{remark}
    Homology is a topic is topology, and specifically algebraic topology,
    where we associate topologucal spaces with algebraic constructions,
    chiefly the abelian group.
    However, in the study of homological algebra,
    we extend these constructions in topology,
    and attempt to gain insight in algbera instead using these same ideas.
\end{remark}

\begin{remark}
    Assume all rings \(R\) are commutative, unless otherwise noted.
\end{remark}


\section{Exact Sequences}\label{sec:exact-sequence}

\begin{definition}
    Suppose \({\{M_i\}}_{i\in\bZ}\) is a sequence of \(R\)-modules,
    with connecting homomorphisms \({\{\phi_i\}}_{i\in\bZ}\), forming
    \begin{equation*}
        \begin{tikzcd}
            \cdots \arrow{r} & M_{i-1} \arrow{r}{\phi_i} &
            M_i \arrow{r}{\phi_{i+1}} & M_{i+1} \arrow{r} & \cdots
        \end{tikzcd}
    \end{equation*}
    We say the sequence is exact at \(M_i\) if \(\img\phi_i = \ker\phi_{i-1}\);
    and the sequence is exact if it is exact at every position.
\end{definition}

\begin{proposition}\label{prop:exact-injective}\label{prop:exact-surjective}
    We have the following two characterizations of homomorphisms:
    \begin{enumerate}[label={(\alph*)}, itemsep=0mm]
        \item \(0 \to M \to N\) exact if and only if \(M \to N\) is injective; and
        \item \(M \to N \to 0\) exact if and only if \(M \to N\) is surjective.
    \end{enumerate}
\end{proposition}
\begin{proof}
    Suppose \(\func{\phi}{M}{N}\) is injective.
    Then \(\ker\phi/\img 0 = \ker\phi = 0\) as desired.
    If it is not injective instead, then \(\ker\phi/0 = \ker\phi \neq 0\) as desired.

    Suppose \(\func{\phi}{M}{N}\) is surjective.
    Then \(\ker 0/\img\phi = N/N = 0\) as desired.
    If it is not surjective, then \(\ker 0/\img\phi \neq N/N = 0\) as desired.
\end{proof}

\begin{definition}
    A short exact sequence is an exact sequence
    \begin{equation*}
        \begin{tikzcd}
            0 \arrow{r} & M_1 \arrow{r} & M_2 \arrow{r} & M_3 \arrow{r} & 0
        \end{tikzcd}
    \end{equation*}
\end{definition}
\begin{proposition}\label{prop:short-exact-sequence}
    Such a sequence is a short exact sequence if and only if \(M_3 \cong M_2/M_1\).
\end{proposition}
\begin{proof}
    Suppose the sequence is exact, with \(\func{f}{M_1}{M_2}\) and \(\func{g}{M_2}{M_3}\).
    We know that \(f\) is an embedding of \(M_1 \subseteq M_2\) by Proposition~\ref{prop:exact-injective}.
    Then \(\ker g/\img f = \ker g/M_1 = 0\), so \(\ker g = M_1\).
    The \hyperref[thm:iso-1-mod]{first isomorphism theorem} then tells us that
    \(M_3 \cong M_2/\ker g = M_2/M_1\).
\end{proof}
\begin{theorem}\label{thm:cut-exact-sequence}
    Every exact sequence can be cut into short exact sequences.
\end{theorem}
\begin{proof}
    Suppose we have the exact sequence
    \begin{equation*}
        M_\bullet: \qquad
        \begin{tikzcd}
            \cdots \arrow{r} & M_{i+1} \arrow{r}{\phi_{i+1}} &
            M_i \arrow{r}{\phi_i} & M_{i-1} \arrow{r} & \cdots
        \end{tikzcd}
    \end{equation*}
    We claim that the following diagram commutes
    \begin{equation*}
        \begin{tikzcd}
            \cdots \arrow{r} & M_{i+1} \arrow{rr}{\phi_{i+1}} \arrow{rd}{\phi_{i+1}} &&
            \textcolor{red}{M_i} \arrow{rr}{\phi_i} \arrow[color=red]{rd}{\phi_i} &&
            M_{i-1} \arrow{r} & \cdots \\
            \cdots && \textcolor{red}{\ker\phi_i} \arrow[color=red]{ru}{\iota_i} \arrow{rd} &&
            \textcolor{red}{\ker\phi_{i-1}} \arrow{ru}{\iota_{i-1}} \arrow[color=red]{rd} && \cdots \\
            \cdots & \textcolor{red}{0} \arrow[color=red]{ru} && 0 \arrow{ru} &&
            \textcolor{red}{0} & \cdots
        \end{tikzcd}
    \end{equation*}
    and the red sequence is exact.
    \(\iota_i\) is the natural inclusion map,
    so clearly \(\phi_{i+1} = \iota_i \circ \phi_{i+1}\).
    By exactness at \(M_i\), we have \(\ker\phi_{i-1} = \img\phi_i\).
    But then the first homomorphism theorem gives us that \(\img\phi_i \cong M_i/\ker\phi_i\),
    which by Proposition~\ref{prop:short-exact-sequence} we know is short exact.
\end{proof}

\begin{lemma}\label{lem:additive-characteristic}
    Suppose \(C\) is a class of \(R\)-modules,
    closed under images and kernels.
    Let \(\func{\lambda}{C}{\bZ}\) be additive.
    If we have a short exact sequence in \(C\),
    \begin{equation*}
        \begin{tikzcd}
            0 \arrow{r} & M_3 \arrow{r} & M_2 \arrow{r} & M_1 \arrow{r} & 0
        \end{tikzcd}
    \end{equation*}
    then \(-\lambda(M_1) + \lambda(M_2) - \lambda(M_3) = 0\).
\end{lemma}
\begin{proof}
    By Proposition~\ref{prop:short-exact-sequence}, we have \(M_1 \cong M_2/M_3\),
    so \(M_2 \cong M_1 \oplus M_3\), hence \(\lambda(M_2) = \lambda(M_1) + \lambda(M_3)\).
\end{proof}
\begin{theorem}\label{thm:additive-characteristic}
    An exact sequence in \(C\) yields \(\sum_{i=1}^n {(-1)}^i \lambda(M_i) = 0\).
\end{theorem}
\begin{proof}
    Let us proceed by induction,
    since the case of \(n=3\) is already proven
    by the \hyperref[lem:additive-characteristic]{lemma above}.
    We shall cut up our sequence using Theorem~\ref{thm:cut-exact-sequence}.
    \begin{equation*}
        \begin{tikzcd}
            \cdots \arrow{r} & M_3 \arrow{rr} \arrow{rd} &&
            \textcolor{red}{M_2} \arrow[color=red]{r}{\phi_2} &
            \textcolor{red}{M_1} \arrow[color=red]{r} & \textcolor{red}{0} \\
            && \textcolor{red}{\ker\phi_2} \arrow[color=red]{ru} \arrow{rd} \\
            & \textcolor{red}{0} \arrow[color=red]{ru} && 0
        \end{tikzcd}
    \end{equation*}
    The red sequence is exact.
    Of course, \(M_3 \to \ker\phi_2 \to 0\) is also exact
    because it is by definition surjective (Proposition~\ref{prop:exact-surjective}).
    Hence by the inductive hypothesis,
    \(\lambda(\ker\phi_2) + \sum_{i=3}^n {(-1)}^i \lambda(M_i) = 0\);
    and in a short exact sequence,
    \(-\lambda(M_1) + \lambda(M_2) - \lambda(\ker\phi_2) = 0\).
    Adding the two equations together yields our desired result.
\end{proof}
\begin{remark}
    The Euler characteristic \(\chi\) obeys, by this same reasoning,
    \begin{equation*}
        \chi(\sigma) = \sum_{i=1}^n {(-1)}^i \dim(V_i) = 0
    \end{equation*}
\end{remark}

\begin{definition}
    Consider a functor \(\func{F}{\textbf{Mod}_R}{\textbf{Mod}_R}\).
    \(F\) is exact if it sends short exact sequences to short exact sequences.
    \(F\) is right-exact if it sends an exact \(M_3 \to M_2 \to M_1 \to 0\)
    to an exact \(F(M_3) \to F(M_2) \to F(M_1) \to 0\).
    \(F\) is left-exact if it sends an exact \(0 \to M_3 \to M_2 \to M_1\)
    to an exact \(0 \to F(M_3) \to F(M_2) \to F(M_1)\).
\end{definition}
\begin{proposition}
    Any exact functor sends exact sequences to exact sequences.
\end{proposition}
\begin{proof}
    % It suffices to prove that any sequence \(M_3 \to M_2 \to M_1\), exact at \(M_2\),
    % gets sent to some sequence \(F(M_3) \to F(M_2) \to F(M_1)\), exact at \(F(M_2)\).
    Suppose we have the exact sequence
    \begin{equation*}
        M_\bullet: \qquad
        \begin{tikzcd}
            \cdots \arrow{r} & M_{i+1} \arrow{r}{\phi_{i+1}} &
            M_i \arrow{r}{\phi_i} & M_{i-1} \arrow{r} & \cdots
        \end{tikzcd}
    \end{equation*}
    and we wish to show that
    \begin{equation*}
        F(M_\bullet): \qquad
        \begin{tikzcd}
            \cdots \arrow{r} & F(M_{i+1}) \arrow{r}{F(\phi_{i+1})} &
            F(M_i) \arrow{r}{F(\phi_i)} & F(M_{i-1}) \arrow{r} & \cdots
        \end{tikzcd}
    \end{equation*}
    is exact at \(F(M_i)\).
    By Theorem~\ref{thm:cut-exact-sequence},
    the following diagram commutes:
    \begin{equation*}
        \begin{tikzcd}
            \cdots \arrow{r} &
            \textcolor{teal}{M_{i+1}} \arrow{rr}{\phi_{i+1}} \arrow[color=teal]{rd}{\phi_{i+1}} &&
            \textcolor{red}{M_i} \arrow{rr}{\phi_i} \arrow[color=red]{rd}{\phi_i} &&
            \textcolor{teal}{M_{i-1}} \arrow{r} & \cdots \\
            \cdots && \textcolor{red}{\ker\phi_i} \arrow[color=red]{ru}{\iota_i} \arrow[color=teal]{rd} &&
            \textcolor{red}{\ker\phi_{i-1}} \arrow[color=teal]{ru}{\iota_{i-1}} \arrow[color=red]{rd} && \cdots \\
            \cdots & \textcolor{red}{0} \arrow[color=red]{ru} && \textcolor{teal}{0} \arrow[color=teal]{ru} &&
            \textcolor{red}{0} & \cdots
        \end{tikzcd}
    \end{equation*}
    and the red and teal sequences are exact.
    Applying the functor \(F\),
    the following diagram commutes:
    \begin{equation*}
        \begin{tikzcd}
            \cdots \arrow{r} &
            \textcolor{teal}{F(M_{i+1})} \arrow{rr}{F(\phi_{i+1})} \arrow[color=teal]{rd}{F(\phi_{i+1})} &&
            \textcolor{red}{F(M_i)} \arrow{rr}{F(\phi_i)} \arrow[color=red]{rd}{F(\phi_i)} &&
            \textcolor{teal}{F(M_{i-1})} \arrow{r} & \cdots \\
            \cdots && \textcolor{red}{F(\ker\phi_i)} \arrow[color=red]{ru}{F(\iota_i)} \arrow[color=teal]{rd} &&
            \textcolor{red}{F(\ker\phi_{i-1})} \arrow[color=teal]{ru}{F(\iota_{i-1})} \arrow[color=red]{rd} && \cdots \\
            \cdots & \textcolor{red}{0} \arrow[color=red]{ru} && \textcolor{teal}{0} \arrow[color=teal]{ru} &&
            \textcolor{red}{0} & \cdots
        \end{tikzcd}
    \end{equation*}
    and the red and teal sequences are once again exact
    (since the two teal sequences are part of a short exact sequence).

    We can now see that \(\img F(\phi_{i+1}) = \img F(\iota_i\circ\phi_{i+1})\).
    But since \(F(\phi_{i+1})\) is surjective onto \(F(\ker\phi_i)\)
    (teal sequence under Proposition~\ref{prop:exact-surjective}),
    \(\img F(\iota_i\circ\phi_{i+1}) = \img F(\iota_i)\).
    By exactness of the red sequence, \(\img F(\iota_i) = \ker F(\phi_i)\).
    As \(F(\iota_{i-1})\) is injective 
    (teal sequence under Proposition~\ref{prop:exact-injective} again),
    \(\ker F(\phi_i) = \ker F(\iota_{i-1}\circ\phi_i)\).
    Lastly, \(\ker F(\iota_{i-1}\circ\phi_i) = \ker F(\phi_i)\) for the black arrow on top.
    Combining these together we have our desired \(\img F(\phi_{i+1}) = \ker F(\phi_i)\).
\end{proof}

\begin{theorem}
    Suppose \(M\) is an \(R\)-module.
    Let us consider two functors
    \(\Hom_R(N,\cdot): N \mapsto \Hom_R(N,M)\) and \(\Hom(\cdot,N): N \mapsto \Hom_R(M,N)\).
    \begin{enumerate}[label={(\alph*)}, itemsep=0mm]
        \item \(0 \to M_3 \to M_2 \to M_1\) is exact
            if and only if \(0 \to \Hom(N,M_3) \to \Hom(N,M_2) \to \Hom(N,M_1)\) is exact; and
        \item \(M_3 \to M_2 \to M_1 \to 0\) is exact
            if and only if \(\Hom(M_3,N) \gets \Hom(M_2,N) \gets \Hom(M_1,N) \gets 0\) is exact.
    \end{enumerate}
    In particular, \(\Hom_R(N,\cdot)\) is a covariant left-exact functor,
    and \(\Hom(\cdot,N)\) is a contravariant right-exact functor.
\end{theorem}
\begin{proof}
    It is easy to see that if \(\func{\phi}{M}{M'}\),
    then under the \(\Hom(N,\cdot)\) functor,
    this maps to \(\func{\phi\circ\cdot}{\Hom(N,M)}{\Hom(N,M')}\),
    as we need to append a map \(M \to M'\);
    this gives us covariance.

    Consider the sequence
    \begin{equation*}
        M_\bullet: \qquad
        \begin{tikzcd}
            0 \arrow{r} & M_3 \arrow{r}{f} & M_2 \arrow{r}{g} & M_1
        \end{tikzcd}
    \end{equation*}
    Suppose \(M_\bullet\) is exact, and let us show that \(\Hom(N,M_\bullet)\) is also exact.
    We first show that \(f\circ\cdot\) is injective.
    If \(\phi\neq\phi'\), then observe that there exists some \(x\) such that \(\phi(x) \neq \phi'(x)\),
    and since \(f\) is injective, \((f\circ\phi)(x) \neq (f\circ\phi')(x)\).
    Hence \(f\circ\phi \neq f\circ\phi'\) as desired.
    Now we show that \(\img(f\circ\cdot) = \ker(g\circ\cdot)\).
    Notice \(\ker(g\circ\cdot) = \{\phi: \img\phi \subseteq \ker g\}\) by definition,
    and \(\img(f\circ\cdot) = \{\phi: \img\phi \subseteq \img f\}\),
    since any function we put into \(f\circ\cdot\) cannot exceed the image,
    and \(f\) is a bijection from \(M_3 \to \img f\),
    so \(f^{-1}\circ\phi\) is the desired map.
    This proves the forward direction of (a).

    Now suppose \(\Hom(N,M_\bullet)\) is exact for all \(N\),
    and show that \(M_\bullet\) is also exact.
    We have the exact sequence
    \begin{equation*}
        \Hom(N,M_\bullet): \qquad
        \begin{tikzcd}
            0 \arrow{r} & \Hom(N,M_3) \arrow{r}{f\circ\cdot} &
            \Hom(N,M_2) \arrow{r}{g\circ\cdot} & \Hom(N,M_1)
        \end{tikzcd}
    \end{equation*}
    Let \(N = \ker f\).
    Consider the embedding \(\func{\iota}{\ker f}{M_3}\), and the zero map.
    We can see that by definition of the kernel, \(f\circ\iota = 0 = f\circ 0\),
    so by injectivity of \(f\circ\cdot\),
    we are forced to conclude that \(\iota = 0\),
    which implies \(\ker f = 0\), and \(f\) is injective.
    We now let \(N = \ker g\), and show that \(\ker g \subseteq \img f\).
    Consider the embedding \(\func{\iota}{\ker g}{M_2}\).
    By definition of the kernel, \(g\circ\iota = 0\),
    so we are forced to conclude that \(\iota \in \ker(g\circ\cdot) = \img(f\circ\cdot)\).
    As we said in the previous paragraph, this gives us \(\img\iota \subseteq \img f\),
    which since \(\iota\) is an injection, \(\ker g \cong \img\iota \subseteq \img f\).
    We finally let \(N = M_3\), and show that \(\img f \subseteq \ker g\).
    Since we know that \(g\circ f\circ\cdot = 0\),
    in particular \(g\circ f\circ\id = g\circ f = 0\).
    This gives us \(\img f \subseteq \ker g\).
    This proves the reverse direction of (a).

    \medskip

    It is easy to see that if \(\func{\phi}{M}{M'}\),
    then under the \(\Hom(\cdot,N)\) functor,
    this maps to \(\func{\cdot\circ\phi}{\Hom(M',N)}{\Hom(M,N)}\),
    as we need to prepend a map \(M \to M'\);
    this gives us contravariance.

    Consider the sequence
    \begin{equation*}
        M_\bullet: \qquad
        \begin{tikzcd}
            M_3 \arrow{r}{f} & M_2 \arrow{r}{g} & M_1 \arrow{r} & 0
        \end{tikzcd}
    \end{equation*}
    Suppose \(M_\bullet\) is exact, and let us show that \(\Hom(M_\bullet,N)\) is also exact.
    We first show that \(\cdot\circ g\) is injective.
    If \(\phi\neq\phi'\), then there exists some \(x\) such that \(\phi(x) \neq \phi'(x)\),
    and since \(g\) is surjective, there exists some \(y\) such that \(g(y) = x\).
    Hence \((\phi\circ g)(y) \neq (\phi'\circ g)(y)\),
    and \(\phi\circ g \neq \phi'\circ g\) as desired.
    Now we show that \(\img(\cdot\circ g) = \ker(\cdot\circ f)\).
    Notice \(\ker(\cdot\circ f) = \{\phi: \img f \subseteq \ker\phi\}\) by definition,
    and \(\img(\cdot\circ g) = \{\phi: \ker g \subseteq \ker\phi\}\),
    since any function \(\psi\circ g\) can only potentially send more elements to zero,
    and the \hyperref[thm:univ-prop-quotient-mod]{universal property of quotients}
    guarantees that \(\phi\) factors through \(g\).
    This proves the forward direction of (b).

    Now suppose \(\Hom(M_\bullet,N)\) is exact for all \(N\),
    and show that \(M_\bullet\) is also exact.
    We have the exact sequence
    \begin{equation*}
        \Hom(M_\bullet,N): \qquad
        \begin{tikzcd}
            \Hom(M_3,N) & \Hom(M_2,N) \arrow{l}[swap]{\cdot\circ f} &
            \Hom(M_1,N) \arrow{l}[swap]{\cdot\circ g} & 0 \arrow{l}
        \end{tikzcd}
    \end{equation*}
    Let \(N = \coker g\), and consider the quotient map \(\func{\pi}{M_1}{\coker g}\).
    Then by definition of the kernel, \(\pi\circ g = 0 = 0\circ g\),
    so by injectivity of \(\cdot\circ g\),
    we are forced to conclude that \(\pi = 0\),
    which implies \(\img g = M_1\), and \(g\) surjective.
    We now let \(N = \coker f\), and show that \(\ker g \subseteq \img f\).
    Consider the projection \(\func{\pi}{M_2}{\coker f}\).
    By definition of the kernel, \(\pi\circ f = 0\),
    so we are forced to conclude that \(\pi \in \ker(\cdot\circ f) = \img(\cdot\circ g)\).
    As we said in the previous paragraph, this gives us \(\ker g \subseteq \ker\pi\),
    and by definition, we have \(\ker g \subseteq \ker\pi = \img f\).
    We finally let \(N = M_1\), and show that \(\img f \subseteq \ker g\).
    Since we know that \(\cdot\circ g\circ f = 0\),
    in particular \(\id\circ g\circ f = g\circ f = 0\).
    This gives us \(\img f \subseteq \ker g\).
    This proves the reverse direction of (b).
\end{proof}

\begin{theorem}[Snake Lemma]\label{thm:snake}\label{lem:snake}
    Suppose we have two exact sequences,
    and the black parts of the following diagram commutes.
    \begin{equation*}
        \begin{tikzcd}
            \textcolor{teal}{0} \arrow[color=teal]{r} & M_1 \arrow{d}{f} \arrow{r} &
            M_2 \arrow{d}{g} \arrow{r} & M_3 \arrow{d}{h} \arrow{r} & 0 \\
            0 \arrow{r} & N_1 \arrow{r} & N_2 \arrow{r} & N_3 \arrow[color=olive]{r} & \textcolor{olive}{0}
        \end{tikzcd}
    \end{equation*}
    Then the black parts of the following sequence is exact.
    \begin{equation*}
        \begin{tikzcd}
            \textcolor{teal}{0} \arrow[color=teal]{r} &
            \ker f \arrow{r} & \ker g \arrow{r} & \ker h \arrow{r}{\delta} &
            \coker f \arrow{r} & \coker g \arrow{r} & \coker h \arrow[color=olive]{r} &
            \textcolor{olive}{0}
        \end{tikzcd}
    \end{equation*}
    Additionally,
    \begin{enumerate}[itemsep=0mm]
        \item[\textcolor{teal}{(a)}] if \(M_1 \to M_2\) is injective, then \(\ker f \to \ker g\) is also injective; and
        \item[\textcolor{olive}{(b)}] if \(N_2 \to N_3\) is surjective, then \(\coker g \to \coker h\) is also surjective.
    \end{enumerate}

    This can be summarized by the following commutative diagram:
    \begin{equation*}
        \begin{tikzcd}
            & 0 \arrow{d} & 0 \arrow{d} & 0 \arrow{d} \\
            \textcolor{teal}{0} \arrow[color=teal]{r} & \textcolor{red}{\ker f} \arrow{d} \arrow[color=red]{r} &
            \textcolor{red}{\ker g} \arrow{d} \arrow[color=red]{r} \arrow[ddd, phantom, ""{coordinate, name=Z}] &
            \textcolor{red}{\ker h} \arrow{d}
            \arrow[rounded corners, dashrightarrow, color=red, to path={
                -- ([xshift=12ex]\tikztostart.east)
                |- (Z) [near end]\tikztonodes
                -| ([xshift=-12ex]\tikztotarget.west)
                -- (\tikztotarget)
            }]{dddll}[at start]{\delta} \\
            \textcolor{teal}{0} \arrow[color=teal]{r} &
            M_1 \arrow[crossing over, "f" description]{d} \arrow{r} &
            M_2 \arrow[crossing over, "g" description]{d} \arrow{r} &
            M_3 \arrow[crossing over, "h" description]{d} \arrow{r} & 0 \\
            0 \arrow{r} & N_1 \arrow{d} \arrow{r} &
            N_2 \arrow{d} \arrow{r} & N_3 \arrow{d} \arrow[color=olive]{r} &
            \textcolor{olive}{0} \\
            & \textcolor{red}{\coker f} \arrow{d} \arrow[color=red]{r} &
            \textcolor{red}{\coker g} \arrow{d} \arrow[color=red]{r} &
            \textcolor{red}{\coker h} \arrow{d} \arrow[color=olive]{r} &
            \textcolor{olive}{0} \\
            & 0 & 0 & 0
        \end{tikzcd}
    \end{equation*}
    where the red sequence is exact, and the teal and olive parts are optional.
\end{theorem}
\begin{proof}
    Let us denote \(\mu\) as the morphisms between \(M_\bullet\),
    and \(\nu\) as the morphisms between \(N_\bullet\).

    We first construct the homomorphism \(\func{\delta}{\ker h}{\coker f}\).
    Consider some \(m_3 \in \ker h \subseteq M_3\).
    Since \(M_2 \to M_3\) is surjective,
    there is some \(m_2 \in M_2\) such that \(\mu(m_2) = m_3\).
    Now let us denote \(n_2 = g(m_2)\).
    We can see that by the diagram,
    \(\nu(n_2) = (\nu\circ g)(m_2) = (h\circ\mu)(m_2) = h(m_3) = 0\).
    Now by exactness, \(n_2 \in \ker(N_2 \to N_3) = \img(N_1 \to N_2)\),
    there exists some \(n_1 \in N_1\) such that \(\nu(n_1) = n_2\).
    We shall then define \(\vfunc{\delta}{\ker h}{\coker f}{m_3}{[n_1]}\),
    where \([n_1]\) denotes the equivalence class of \(n_1\) in \(\coker f = N_1/\img f\).
    Let us now check that \(\delta\) is indeed well-defined.
    Suppose we have chosen another \(m_2' \in M_2\) such that \(\mu(m_2') = m_3\).
    Then exactness yields \(m_2-m_2' \in \ker(M_2 \to M_3) = \img(M_1 \to M_2)\),
    where we can say there exists some \(\tilde{m}_1 \in M_1\) such that \(\mu(\tilde{m}_1) = m_2-m_2'\).
    Let us now define \(n_2' = g(m_2')\).
    Observe that \(\nu(n_2-n_2') = (\nu\circ g)(m_2-m_2') = (h\circ\mu)(m_2-m_2') = h(0) = 0\).
    By exactness, there is some \(\tilde{n}_1 \in N_1\) such that \(\nu(\tilde{n}_1) = n_2-n_2'\).
    But then we see that \(\nu(\tilde{n}_1-f(\tilde{m}_1)) = 0 \in N_2\)
    since they both map into the same object,
    so by injectivity, \(f(\tilde{m}_1) = \tilde{n}_1\).
    Lastly, we can see that by exactness, \([\tilde{n}_1] = [f(\tilde{m}_1)] = [0]\),
    so the choice of elements do not affect our definition.

    We then inherit the maps \(\ker f \to \ker g \to \ker h\) from \(M_1 \to M_2 \to M_3\).
    Let us quickly check that this is a valid definition.
    Suppose \(m_1 \in \ker f \subseteq M_1\).
    We claim that \(m_2 = \mu(m_1) \in \ker g \subseteq M_2\).
    We have \(f(m_1) = 0\), so \(g(m_2) = (g\circ\mu)(m_1) = (\nu\circ f)(m_1) = \nu(0) = 0\) as desired.
    This validates \(\ker f \to \ker g\).
    To validate \(\ker g \to \ker h\),
    simply substitute \(M_1 \mapsto M_2\), \(M_2 \mapsto M_3\), \(f \mapsto g\), and \(g \mapsto h\).

    We also wish to define \(\coker f \to \coker g \to \coker h\)
    by sending the coset representatives to coset representatives in \(N_1 \to N_2 \to N_3\);
    i.e.\ \(\bar{\nu}[n] = [\nu(n)]\).
    Let us quickly check that this is also a valid definition.
    Suppose \(n_1,n_1' \in N_1\) such that \(n_1' \in [n_1]\).
    Let \(\nu(n_1) = n_2\) and \(\nu(n_1') = n_2'\).
    Then we have \(n_1-n_1' \in [0] = \img f\),
    so there exists some \(\tilde{m}_1 \in M_1\) such that \(f(\tilde{m}_1) = n_1-n_1'\).
    But the diagram tells us that
    \((g\circ\mu)(\tilde{m}_1) = (\nu\circ f)(\tilde{m}_1) = \nu(n_1-n_1') = n_2-n_2'\),
    so \(n_2-n_2' \in \img g = [0]\) indeed.
    This validates \(\func{\bar{\nu}}{\coker f}{\coker g}\).
    By similarity as seen in the other diagram,
    this also validates \(\func{\bar{\nu}}{\coker g}{\coker h}\).

    We prove exactness at \(\ker f\);
    that is, \(\func{\mu}{\ker f}{\ker g}\) is injective,
    assuming \(\func{\mu}{M_1}{M_2}\) is injective.
    This is obvious, since the homomorphism is inherited.

    We prove exactness at \(\ker g\);
    that is, \(\img(\ker f \to \ker g) = \ker(\ker g \to \ker h)\).
    Firstly, the forward inclusion is obvious,
    since for any \(m_1 \in \ker f \subseteq M_1\),
    exactness in \(M_\bullet\) gives \(\mu^2(m_1) = 0\).
    Secondly, for the reverse inclusion,
    suppose we have an element \(m_2 \in \ker g \subseteq M_2\) such that \(\mu(m_2) = 0\).
    By exactness of \(M_\bullet\), there exists \(m_1 \in M_1\) such that \(\mu(m_1) = m_2\).
    But since \((\nu\circ f)(m_1) = (g\circ\mu)(m_1) = g(m_2) = 0\),
    by injectivity of \(N_1 \to N_2\), \(f(m_1) = 0\).
    Hence \(m_1 \in \ker f\), and \(m_2 \in \img(\ker f \to \ker g)\).

    We prove exactness at \(\ker h\);
    that is, \(\img(\ker g \to \ker h) = \ker\delta\).
    Firstly, for the forward inclusion,
    suppose we have \(m_2 \in \ker g \subseteq M_2\) and \(m_3 \in \ker h \subseteq M_3\)
    such that \(\mu(m_2) = m_3\).
    Then we see that \(n_2 = g(m_2) = 0\),
    so by injectivity of \(N_1 \to N_2\),
    the \(n_1 \in N_1\) such that \(\nu(n_1) = n_2 = 0\) implies \(n_1 = 0\),
    so \(\delta(m_3) = 0 \in \coker f\).
    Secondly, for the reverse inclusion,
    suppose \([n_1] = \delta(m_3) = 0\) for some \(m_3 \in \ker h \subseteq M_3\).
    Let us denote \(m_2\) and \(n_2\) as in the definition of \(\delta\),
    and note that in particular \(g(m_2) = n_2\).
    But now, we have that \(n_1 \in \img f\) so there exists some \(m_1 \in M_1\)
    such that \(f(m_1) = n_1\).
    We can define \(m_2' = \mu(m_1) \in M_2\), and see that \(\mu^2(m_1) = 0\).
    In particular, we have \(g(m_2) = (g\circ\mu)(m_1) = (\nu\circ f)(m_1) = \nu(n_1) = n_2\).
    We can now see that \(\mu(m_2-m_2') = m_3\),
    and it suffices to prove that \(m_2-m_2' \in \ker g\).
    This is given since \(g(m_2-m_2') = n_2 - n_2 = 0\).
    
    We prove exactness at \(\coker f\);
    that is, \(\img\delta = \ker(\coker f \to \coker g)\).
    Firstly, for the forward inclusion,
    suppose we have \(n_1 \in N_1\)
    such that there exists \(m_3 \in \ker h\) with \(\delta(m_3) = [n_1]\).
    Let us denote \(m_2\) and \(n_2\) as in the definition of \(\delta\),
    and in particular, \(g(m_2) = n_2\).
    This exactly tells us that \(n_2 \in \img g\),
    so \(n_2 \in [0]\).
    Now as \(\nu(n_1) = n_2\), we have \(\bar{\nu}[n_1] = [n_2] = [0]\) as desired.
    Secondly, for the reverse inclusion,
    suppose \(\bar{\nu}[n_1] = [0]\) for some \(n_1 \in N_1\).
    Then \(\nu(n_1) = n_2 \in \img g\),
    so there exists \(m_2 \in M_2\) such that \(g(m_2) = n_2\).
    Let us denote \(\mu(m_2) = m_3 \in M_3\).
    We claim that this \(\delta(m_3) = [n_1]\).
    Observe that in the definition of \(\delta\) we proved that
    we can choose any \(m_2 \in M_2\) such that \(\mu(m_2) = m_3\),
    and the corresponding \([n_1]\) that we get will be the same.
    Hence the \([n_1]\) that we started with is the correct one.

    We prove exactness at \(\coker g\);
    that is, \(\img(\coker f \to \coker g) = \ker(\coker g \to \coker h)\).
    Firstly, the forward inclusion is obvious,
    since for any \([m_1] \in \coker f\),
    \(\bar{\nu}^2[m_1] = [\nu^2(m_1)] = [0] \in \coker h\)
    by exactness in \(N_\bullet\).
    Secondly, for the reverse inclusion,
    suppose we have an element \([n_2] \in \coker g\) such that \(\bar{\nu}[n_2] = [0]\).
    Consider \(\nu(n_2) = n_3 \in N_3\), and in particular \(n_3 \in [0]\).
    Then we know that \(n_3 \in \img h\),
    so there exists \(m_3 \in M_3\) such that \(h(m_3) = n_3\).
    By surjectivity of \(M_2 \to M_3\),
    there exists \(m_2 \in M_2\) such that \(\mu(m_2) = m_3\).
    Let us define \(n_2' = g(m_2)\),
    and see that \(\nu(n_2-n_2') = n_3-(\nu\circ g)(m_2) = n_3-(h\circ\mu)(m_2) = n_3 - h(m_3) = n_3-n_3 = 0\).
    Hence by exactness of \(N_\bullet\),
    there is \(\tilde{n}_1 \in N_1\) such that \(\nu(\tilde{n}_1) = n_2-n_2'\).
    But observe that \(n_2' \in \img g\), so \([n_2-n_2'] = [n_2]-[n_2'] = [n_2]-[0] = [n_2]\),
    and we have \(\bar{\nu}[\tilde{n}_1] = [\nu(\tilde{n}_1)] = [n_2-n_2'] = [n_2]\).

    We prove exactness at \(\coker h\);
    that is, \(\func{\bar{\nu}}{\coker g}{\coker h}\) is surjective,
    assuming \(\func{\nu}{N_2}{N_3}\) is surjective.
    This is obvious, since if \(n_3 = \nu(n_2)\),
    then any \([n_3] = [\nu(n_2)] = \bar{\nu}[n_2]\).1
\end{proof}

\begin{corollary}
    Suppose we have two short exact sequences,
    and the following diagram commutes.
    \begin{equation*}
        \begin{tikzcd}
            0 \arrow{r} & M_1 \arrow{d}{f} \arrow{r} &
            M_2 \arrow{d}{g} \arrow{r} & M_3 \arrow{d}{h} \arrow{r} & 0 \\
            0 \arrow{r} & N_1 \arrow{r} & N_2 \arrow{r} & N_3 \arrow{r} & 0
        \end{tikzcd}
    \end{equation*}
    Then if any two of \(f,g,h\) are isomorphisms,
    then the third one must also be an isomorphism.
\end{corollary}
\begin{proof}
    Suppose \(f,g\) are isomorphisms.
    Then \(\ker f = \ker g = \coker f = \coker g = 0\),
    and by the \hyperref[lem:snake]{snake lemma}, we have the exact sequence
    \begin{equation*}
        \begin{tikzcd}
            0 \arrow{r} & 0 \arrow{r} & 0 \arrow{r} & \ker h \arrow{r}{\delta} &
            0 \arrow{r} & 0 \arrow{r} & \coker h \arrow{r} & 0
        \end{tikzcd}
    \end{equation*}
    Hence clearly we have \(\ker h = \coker h = 0\).
    
    Suppose \(f,h\) are isomorphisms.
    Then \(\ker f = \ker h = \coker f = \coker h = 0\),
    and by the \hyperref[lem:snake]{snake lemma}, we have the exact sequence
    \begin{equation*}
        \begin{tikzcd}
            0 \arrow{r} & 0 \arrow{r} & \ker g \arrow{r} & 0 \arrow{r}{\delta} &
            0 \arrow{r} & \coker g \arrow{r} & 0 \arrow{r} & 0
        \end{tikzcd}
    \end{equation*}
    Hence clearly we have \(\ker g = \coker g = 0\).
    
    Suppose \(g,h\) are isomorphisms.
    Then \(\ker g = \ker h = \coker g = \coker h = 0\),
    and by the \hyperref[lem:snake]{snake lemma}, we have the exact sequence
    \begin{equation*}
        \begin{tikzcd}
            0 \arrow{r} & \ker f \arrow{r} & 0 \arrow{r} & 0 \arrow{r}{\delta} &
            \coker f \arrow{r} & 0 \arrow{r} & 0 \arrow{r} & 0
        \end{tikzcd}
    \end{equation*}
    Hence clearly we have \(\ker f = \coker f = 0\).
\end{proof}

\begin{lemma}[Four Lemmas]\label{lem:four}
    Suppose we have two exact sequences of four objects each,
    and homomorphisms \(f_\bullet\) between the two sequences.
    \begin{enumerate}[label={(\alph*)}, itemsep=0mm]
        \item If \(f_2,f_4\) are injective, and \(f_1\) is surjective,
            then \(f_3\) is injective.
            \begin{equation*}
                \begin{tikzcd}
                    & 0 \arrow{d} & 0 \arrow[dashrightarrow]{d} & 0 \arrow{d} \\
                    M_1 \arrow{r} \arrow{d}{f_1} & M_2 \arrow{r} \arrow{d}{f_2} &
                    M_3 \arrow{r} \arrow{d}{f_3} & M_4 \arrow{d}{f_4} \\
                    N_1 \arrow{r} \arrow{d} & N_2 \arrow{r} & N_3 \arrow{r} & N_4 \\
                    0
                \end{tikzcd}
            \end{equation*}
        \item If \(f_2,f_4\) are surjective, and \(f_5\) is injective,
            then \(f_3\) is surjective.
            \begin{equation*}
                \begin{tikzcd}
                    & & & 0 \arrow{d} \\
                    M_2 \arrow{r} \arrow{d}{f_2} & M_3 \arrow{r} \arrow{d}{f_3} &
                    M_4 \arrow{d}{f_4} \arrow{r} & M_5 \arrow{d}{f_5} \\
                    N_2 \arrow{r} \arrow{d} & N_3 \arrow{r} \arrow[dashrightarrow]{d} &
                    N_4 \arrow{r} \arrow{d} & N_5 \\
                    0 & 0 & 0
                \end{tikzcd}
            \end{equation*}
    \end{enumerate}
\end{lemma}
\begin{proof}
    Let us denote \(\mu\) as the morphisms between \(M_\bullet\),
    and \(\nu\) as the morphisms between \(N_\bullet\).

    For (a), we wish to prove that given some element \(m_3 \in M_3\),
    if \(f(m_3) = 0\), then \(m_3 = 0\).
    Consider such an \(m_3 \in M_3\).
    Then we have \((f\circ\mu)(m_3) = (\nu\circ f)(m_3) = \nu(0) = 0\).
    But as \(f_4\) is injective, \(\mu(m_3) = 0\),
    so by exactness of \(M_\bullet\),
    there exists some \(m_2 \in M_2\) such that \(\mu(m_2) = m_3\).
    Denote \(n_2 = f(m_2)\).
    Now we have \(\nu(n_2) = (\nu\circ f)(m_2) = (f\circ\mu)(m_2) = f(m_3) = 0\),
    and hence by exactness of \(N_\bullet\),
    there exists \(n_1 \in N_1\) such that \(\nu(n_1) = n_2\).
    Since \(f_1\) is surjective, we can find \(m_1 \in M_1\) such that \(f(m_1) = n_1\).
    Observe that \((f\circ\mu)(m_1) = (\nu\circ f)(m_1) = \nu(n_1) = n_2 = f(m_2)\),
    so due to injectivity of \(f_2\), \(\mu(m_1) = m_2\).
    Now by exactness of \(M_\bullet\), \(m_3 = \mu(m_2) = \mu^2(m_1) = 0\) as desired.

    For (b), we wish to prove that given any \(n_3 \in N_3\),
    there exists \(m_3 \in M_3\) such that \(f(m_3) = n_3\).
    Let \(n_4 = \nu(n_3)\).
    By surjectivity of \(f_4\), there exists \(m_4 \in M_4\) such that \(f(m_4) = n_4\).
    Due to exactness of \(N_\bullet\), we have \(\nu(n_4) = \nu^2(n_3) = 0\).
    Then \((f\circ\mu)(m_4) = (\nu\circ f)(m_4) = \nu(n_4) = 0\),
    so as \(f_5\) injective, \(\mu(m_4) = 0\).
    Since \(M_\bullet\) is exact, there exists \(m_3' \in M_3\) such that \(\mu(m_3') = m_4\).
    Allow us to denote \(f(m_3') = n_3'\).
    Since we have \(\nu(n_3-n_3') = n_4 - (\nu\circ f)(m_3') = n_4 - (f\circ\mu)(m_3') = n_4 - f(m_4) = n_4-n_4 = 0\),
    by exactness in \(N_\bullet\) again,
    there exists some \(\tilde{n}_2 \in N_2\) such that \(\nu(\tilde{n}_2) = n_3-n_3'\).
    By surjectivity of \(f_2\), we have \(\tilde{m}_2 \in M_2\) such that \(f(\tilde{m}_2) = \tilde{n}_2\).
    Let \(\tilde{m}_3 = \mu(\tilde{m}_2)\).
    We have \(f(\tilde{m}_3) = (f\circ\mu)(\tilde{m}_2) = (\nu\circ f)(\tilde{m}_2) = \nu(\tilde{n}_2) = n_3-n_3'\),
    so there exists some \(m_3 = \tilde{m}_3 + m_3' \in M_3\)
    such that \(f(m_3) = n_3-n_3'+n_3' = n_3\).
\end{proof}

\begin{corollary}[Five Lemma]\label{cor:five}\label{lem:five}
    Suppose we have two exact sequences of five objects each,
    and homomorphisms \(f_\bullet\) between the two sequences.
    If \(f_1\) is surjective, \(f_2,f_4\) are isomorphisms, and \(f_5\) is injective,
    then \(f_3\) is an isomorphism.
    \begin{equation*}
        \begin{tikzcd}
            & 0 \arrow{d} & 0 \arrow[dashrightarrow]{d} & 0 \arrow{d} & 0 \arrow{d} \\
            M_1 \arrow{r} \arrow{d}{f_1} & M_2 \arrow{r} \arrow{d}{f_2} &
            M_3 \arrow{r} \arrow{d}{f_3} & M_4 \arrow{r} \arrow{d}{f_4} & M_5 \arrow{d}{f_5} \\
            N_1 \arrow{r} \arrow{d} & N_2 \arrow{r} \arrow{d} &
            N_3 \arrow{r} \arrow[dashrightarrow]{d} & N_4 \arrow{r} \arrow{d} & N_5 \\
            0 & 0 & 0 & 0
        \end{tikzcd}
    \end{equation*}
\end{corollary}
\begin{proof}
    Simply combine the diagrams of the two \hyperref[lem:four]{four lemmas}.
\end{proof}


\section{Chain Complexes}

\begin{definition}
    Suppose \({\{C_i\}}_{i\in\bZ}\) are \(R\)-modules.
    A chain complex \((C_\bullet,d_\bullet)\) is a chain of modules
    alongside their connecting homomorphisms
    \begin{equation*}
        C_\bullet: \qquad
        \begin{tikzcd}
            \cdots \arrow{r} & C_{i+1} \arrow{r}{d_{i+1}} &
            C_i \arrow{r}{d_i} & C_{i-1} \arrow{r} & \cdots
        \end{tikzcd}
    \end{equation*}
    where \(d^2 = d_n \circ d_{n+1} = 0\).
    The \(i\)-chains are \(C_i\),
    the \(i\)-cycles are \(\ker d_i \subseteq C_i\),
    and the \(i\)-boundaries are \(\img d_{i+1} \subseteq C_i\).
\end{definition}
\begin{definition}
    Suppose \((C_\bullet,d_\bullet)\) is a chain complex.
    The \(i\)th homology is \(H_i(C_\bullet) = \ker d_i/\img d_{i+1}\).
    The chain complex is acyclic or an exact sequence if \(H_i = 0\) for all \(i\).
\end{definition}

\begin{definition}
    The dual to a chain complex is a cochain complex \((C^\bullet,d^\bullet)\)
    \begin{equation*}
        C^\bullet: \qquad
        \begin{tikzcd}
            \cdots \arrow{r} & C^{i-1} \arrow{r}{d^{i-1}} &
            C^i \arrow{r}{d^i} & C^{i+1} \arrow{r} & \cdots
        \end{tikzcd}
    \end{equation*}
    where \(d^2 = d^n \circ d^{n-1} = 0\).
    The \(i\)-cochains are \(C^i\),
    the \(i\)-cocycles are \(\ker d^i \subseteq C^i\),
    and the \(i\)-coboundaries are \(\img d^{i-1} \subseteq C^i\).
\end{definition}
\begin{definition}
    Suppose \((C^\bullet,d^\bullet)\) is a cochain complex.
    The \(i\)th cohomology is \(H^i(C^\bullet) = \ker d^i/\img d^{i-1}\).
\end{definition}

\begin{definition}
    Suppose \(C_\bullet,D_\bullet\) are chain complexes of \(R\)-modules.
    A map of complexes \(\func{f_\bullet}{C_\bullet}{D_\bullet}\)
    are a set of module homomorphisms \(\func{f_i}{C_i}{D_i}\)
    such that the following diagram commutes.
    \begin{equation*}
        \begin{tikzcd}
            \cdots \arrow{r} & C_{i+1} \arrow{r} \arrow{d}{f_{i+1}} &
            C_i \arrow{r} \arrow{d}{f_i} & C_{i-1} \arrow{r} \arrow{d}{f_{i-1}} & \cdots \\
            \cdots \arrow{r} & D_{i+1} \arrow{r} & D_i \arrow{r} & D_{i-1} \arrow{r} & \cdots
        \end{tikzcd}
    \end{equation*}
\end{definition}

\begin{proposition}
    The category \(\textbf{Ch}_R\) consisting of all chain complexes of \(R\)-modules \(C_\bullet\)
    and map of complexes \(\func{f_\bullet}{C_\bullet}{D_\bullet}\) forms an abelian category.
\end{proposition}
\begin{proof}
    Composition and association of homomorphisms are obvious,
    since they are inherited from module homomorphisms.
    The identity morphisms is simply \(\func{\id_\bullet}{C_\bullet}{C_\bullet}\)
    with \(\func{\id_i}{C_i}{C_i}\) being the usual identity.

    Then abelianness of the category is inherited from
    \hyperref[prop:module-ab-cat]{the abelianness of \(\textbf{Mod}_R\)}.
    The kernel and cokernel complexes are then constructed the same way
    as in the proof of the \hyperref[lem:snake]{snake lemma}.
\end{proof}

\begin{proposition}
    The \(i\)th homology is a functor \(\func{H_i}{\textbf{Ch}_R}{\textbf{Mod}_R}\).
\end{proposition}
\begin{proof}
    We have already defined the map between objects,
    so it suffices to define \(\func{H_i(f_\bullet)}{H_i(C_\bullet)}{H_i(D_\bullet)}\).
    Consider the following part of a commutative diagram.
    \begin{equation*}
        \begin{tikzcd}
            \cdots \arrow{r} & C_{i+1} \arrow{r} \arrow{d}{f_{i+1}} &
            C_i \arrow{r} \arrow{d}{f_i} & C_{i-1} \arrow{r} \arrow{d}{f_{i-1}} & \cdots \\
            \cdots \arrow{r} & D_{i+1} \arrow{r} & D_i \arrow{r} & D_{i-1} \arrow{r} & \cdots
        \end{tikzcd}
    \end{equation*}

    We first show that \(f\) maps kernels to kernels.
    Suppose \(c_i \in \ker d_i\).
    Then \((d \circ f)(c_i) = (f \circ d)(c_i) = f(0) = 0 \in D_{i-1}\).
    Hence \(f(c_i) \in \ker d_i\) too.

    We then show that \(f\) maps images to images.
    Suppose \(c_i \in \img d_{i+1}\).
    Then there exists \(c_{i+1} \in C_{i+1}\) such that \(d(c_{i+1}) = c_i\).
    But we have \((d \circ f)(c_{i+1}) = (f \circ d)(c_{i+1}) = f(c_i)\),
    so \(f(c_i) \in \img d_{i+1}\) too.

    We can then show that \(H_i(f_\bullet)[c_i] = [f(c_i)]\) is a valid definition.
    But we have exactly proven that \(c_i \in [0] = \img d_{i+1}\)
    implies \(f(c_i) \in [0] = \img d_{i+1}\) as desired.
\end{proof}
\begin{definition}
    \(\func{f_\bullet}{C_\bullet}{D_\bullet}\) is a quasi-isomorphism
    if all of the homology functors \(\func{H_i(f_\bullet)}{H_i(C_\bullet)}{H_i(D_\bullet)}\) are isomorphisms.
\end{definition}

\begin{definition}
    Suppose \(M\) is an \(R\)-module.
    A resolution of \(M\) is an exact sequence
    \begin{equation*}
        \begin{tikzcd}
            \cdots \arrow{r} & C_2 \arrow{r} & C_1 \arrow{r} & C_0 \arrow{r} &
            M \arrow{r} & 0
        \end{tikzcd}
    \end{equation*}
    For any property \(P\) of modules,
    a \(P\) resolution is a resolution such that all \(C_i\) have this property \(P\).
\end{definition}
\begin{lemma}
    Every module has a free resolution.
\end{lemma}
\begin{proof}
    By the \hyperref[ax:choice]{axiom of choice},
    suppose we have generators \({\{m_j\}}_{j \in I} \subset M\).
    Then we can set \(C_0 = \bigoplus m_j R\).
    Inductively, we can generate \({\{k_j\}}_{j \in I} \subset \ker C_i\),
    and then set \(C_{j+1} = \bigoplus k_j R\).
\end{proof}


\section{Derived Functors}


\section{Group Cohomology}
